%%% FIELD prop-known-noise-statement
Suppose that the common measurement covariance \(E\) is supplied, and set \(\widetilde C_k=C_k-E\).  Within the present stable, principal-log-embeddable, Gaussian, zero-process-noise structural model, the transition fiber compatible with the four observed marginal laws is exactly \(\mathcal F_{\mathrm{known}}(E)\).  Its covariance restrictions are \(W\widetilde C_kW^{\mathsf T}=\widetilde C_{k+1}\) for \(k=0,1,2\), and its residual-noise restriction is the equality \(E_W=E\), equivalently \(C_0-E-T_W=0\), rather than only the unknown-noise inequality \(C_0-T_W\succ0\).  With only snapshots \(k=0,1,2\), the covariance and affine-mean consequences are exactly \(W\widetilde C_0W^{\mathsf T}=\widetilde C_1\), \(W\widetilde C_1W^{\mathsf T}=\widetilde C_2\), and \(Wd_0=d_1\).  Under the identification of the published discrete transition with \(W\), its state covariances with \(\widetilde C_k\), and its known process covariance with \(Q=0\), these are an equation-level subclass of AaltoLamolineGoncalves2022, Theorem 2.1 and Corollary 2.1.  The inclusion is strict at model scope: their full-state system permits a known nonzero process covariance, an invertible discrete-time transition without stability or continuous-time embeddability restrictions, and finite-second-moment non-Gaussian laws, whereas it has no measurement-error layer.  Thus only the displayed moment-equation consequence, not equality of observation-model classes, is inherited.  If the two three-snapshot whitened congruence invariants have the same simple spectrum and every source probe coordinate is nonzero, the three-snapshot equations are computably reduced to at most \(2^p\) sign candidates.  The candidates for \(\mathcal F_{\mathrm{known}}(E)\) are then tested against the remaining fourth-snapshot congruence and probe equation and against the stability, principal-logarithm, Stein, and equality \(E_W=E\) filters.
%%% FIELD aalto-model-scope-statement
Section 2 of AaltoLamolineGoncalves2022 studies, in a locally denoted dimension \(n\), the exactly observed full-state recursion \(Z_j=A Z_{j-1}+b+\eta_{j-1}\), where the process innovations have known covariance \(Q\) and the temporal orthogonality specified there.  Its moment equations are \(\bar m_j=A\bar m_{j-1}+b\) and \(P_j=AP_{j-1}A^{\mathsf T}+Q\).  Theorem 2.1 assumes that \(A\) is invertible but imposes neither stability nor continuous-time embeddability, and the paragraph immediately following its proof states that finite second moments, rather than Gaussianity, suffice for these moment equations and the stated identification argument.
%%% FIELD prop-known-noise-proof
Let \(\mathcal C_4(E)\) be the set of transition matrices occurring in structural explanations of the four observed marginal laws with the supplied covariance \(E\), subject to \(\text{ass:spacing}\), \(\text{ass:linear-flow}\), \(\text{ass:hurwitz}\), \(\text{ass:alias-strip}\), \(\text{ass:initial-gaussian}\), \(\text{ass:common-measurement-noise}\), \(\text{ass:measurement-independence}\), and \(\text{ass:no-process-noise}\).  This is proof notation for a causal-model compatibility set; it is not a definition of the observable identifying functional.  We first prove
\[
\mathcal C_4(E)=\mathcal F_{\mathrm{known}}(E).
\]

For necessity, take \(W\in\mathcal C_4(E)\), and write \(L_k=\operatorname{Cov}(X_k)\).  By \(\text{ass:linear-flow}\), \(\text{ass:no-process-noise}\), and \(\text{ass:initial-gaussian}\),
\[
 \mathbb E[X_k]=W^k(\mu_0-x^\star)+x^\star,
 \qquad
 L_k=W^k\Sigma(W^{\mathsf T})^k.
\]
Because \(W=\exp(hG)\) is invertible and \(\Sigma\succ0\), congruence by \(W^k\) gives \(L_k\succ0\) for every \(k\in K\).  The errors have mean zero by \(\text{ass:common-measurement-noise}\), and their independence from \(X_0\), hence from each deterministic \(X_k\), follows from \(\text{ass:measurement-independence}\).  Therefore
\[
 m_k=\mathbb E[X_k],
 \qquad
 C_k=L_k+E,
 \qquad
 \widetilde C_k=C_k-E=L_k\succ0.
\]
The deterministic covariance recursion and the affine mean formula now yield, for every index for which the displayed objects exist,
\[
 W\widetilde C_kW^{\mathsf T}=\widetilde C_{k+1}\quad(k=0,1,2),
 \qquad
 d_k=W^k(W-I_p)(\mu_0-x^\star),
\]
and hence
\[
 Wd_0=d_1,
 \qquad
 Wd_1=d_2.
\]

We next verify the filters in \(\text{def:known-noise-fiber}\).  On a complex Jordan block of \(hG\) with eigenvalue \(hg\), exponentiation produces the eigenvalue \(e^{hg}\).  Thus \(\text{ass:hurwitz}\) and \(h>0\) from \(\text{ass:spacing}\) give
\[
 \rho(W)=\max_{g\in\operatorname{spec}(G)}e^{h\operatorname{Re}(g)}<1.
\]
If \(e^{hg}\) were a negative real number, then \(h\operatorname{Im}(g)\) would equal \((2\ell+1)\pi\) for some integer \(\ell\), which is impossible under \(\text{ass:alias-strip}\); exponentials are nonzero.  Hence
\[
 \operatorname{spec}(W)\cap(-\infty,0]=\varnothing.
\]
The common covariance cancels from the first adjacent difference:
\[
 H_0=C_0-C_1=L_0-L_1=L_0-WL_0W^{\mathsf T}.
\]
If \(D-WDW^{\mathsf T}=0\), vectorization gives
\[
 \bigl(I_{p^2}-W\otimes W\bigr)\operatorname{vec}(D)=0.
\]
Every eigenvalue of \(W\otimes W\) is a product of two eigenvalues of \(W\), so its spectral radius is \(\rho(W)^2<1\).  Consequently \(I_{p^2}-W\otimes W\) is invertible, and the Stein equation has at most one solution.  Since \(L_0\) is a solution, the reconstruction in the frozen definition satisfies
\[
 T_W=L_0=\widetilde C_0,
 \qquad
 E_W=C_0-T_W=C_0-\widetilde C_0=E.
\]
All restrictions in \(\text{def:known-noise-fiber}\) hold, proving \(\mathcal C_4(E)\subseteq\mathcal F_{\mathrm{known}}(E)\).

For sufficiency, take \(W\in\mathcal F_{\mathrm{known}}(E)\), and define
\[
 \Sigma=widetilde C_0,
 \qquad
 \mu_0=m_0,
 \qquad
 b=m_1-Wm_0.
\]
Membership in the fiber gives \(\Sigma\succ0\) and \(\rho(W)<1\).  In particular, \(1\notin\operatorname{spec}(W)\), so \(I_p-W\) is invertible and
\[
 x^\star=(I_p-W)^{-1}b
\]
is well defined.  Since \(b=(I_p-W)x^\star\), the affine recursion beginning at \(m_0\) gives \(m_1\) at its first step.  The two probe restrictions then give successively
\[
 \begin{aligned}
 Wm_1+b
 &=m_1+W(m_1-m_0)
 =m_1+Wd_0
 =m_1+d_1=m_2,\\
 Wm_2+b
 &=m_2+W(m_2-m_1)
 =m_2+Wd_1
 =m_2+d_2=m_3.
 \end{aligned}
\]

The spectral-cut condition permits the principal logarithm.  More explicitly, on a complex Jordan block \(J=\lambda I+N\), where \(N^s=0\) and \(-\pi<\operatorname{Arg}(\lambda)<\pi\), set
\[
 \operatorname{Log}(J)
 =\bigl(\log|\lambda|+i\operatorname{Arg}(\lambda)\bigr)I
 +\sum_{r=1}^{s-1}\frac{(-1)^{r+1}}{r}\left(\frac{N}{\lambda}\right)^r.
\]
The finite nilpotent logarithm and exponential series are inverse under composition, so exponentiating this block returns \(J\).  Conjugate blocks receive conjugate logarithms; pairing them in a real canonical form shows that \(\operatorname{Log}(W)\) is real.  Its eigenvalues are the scalar principal logarithms of the eigenvalues of \(W\).  Therefore
\[
 G_W=h^{-1}\operatorname{Log}(W)
\]
is real and satisfies
\[
 \exp(hG_W)=W,
 \qquad
 \operatorname{Re}\operatorname{spec}(G_W)
 =h^{-1}\log|\operatorname{spec}(W)|<0,
 \qquad
 -\frac{\pi}{h}<\operatorname{Im}\operatorname{spec}(G_W)<\frac{\pi}{h}.
\]
Thus it satisfies \(\text{ass:hurwitz}\) and \(\text{ass:alias-strip}\).  Conversely, if \(L\) is any real logarithm of \(W\) whose spectrum lies in the open strip \(-\pi<\operatorname{Im}(z)<\pi\), the scalar identity \(\operatorname{Log}(e^z)=z\) is holomorphic on a neighborhood of \(\operatorname{spec}(L)\).  Composition of the corresponding convergent local power series on each generalized eigenspace gives
\[
 L=\operatorname{Log}(e^L)=\operatorname{Log}(W).
\]
Hence the open alias strip leaves no additional generator branch.

Construct \(X_0\sim\mathcal N(m_0,\Sigma)\), set \(X_k=W^k(X_0-x^\star)+x^\star\), and take mutually independent \(\varepsilon_k\sim\mathcal N(0,E)\), independent of \(X_0\).  This construction satisfies the stated Gaussian, independence, and zero-process-noise assumptions.  The preceding recursion gives \(\mathbb E[X_k]=m_k\).  The level congruences in \(\text{def:known-noise-fiber}\) give inductively
\[
 \operatorname{Cov}(X_k)
 =W^k\Sigma(W^{\mathsf T})^k
 =\widetilde C_k
 \qquad(k\in K).
\]
Consequently
\[
 \mathbb E[X_k+\varepsilon_k]=m_k,
 \qquad
 \operatorname{Cov}(X_k+\varepsilon_k)=\widetilde C_k+E=C_k.
\]
Both the constructed and observed marginal laws are nonsingular Gaussian laws, so equality of these moments is equality of laws.  Thus \(W\in\mathcal C_4(E)\), proving the reverse inclusion.  Moreover \(E_W=E\succ0\) by \(\text{ass:common-measurement-noise}\); fixing this residual is strictly stronger than merely requiring \(C_0-T_W\succ0\).

Now retain only snapshots \(0,1,2\).  The same calculations show that their covariance and mean consequences are precisely
\[
 W\widetilde C_0W^{\mathsf T}=\widetilde C_1,
 \qquad
 W\widetilde C_1W^{\mathsf T}=\widetilde C_2,
 \qquad
 Wd_0=d_1.
\]
Conversely, these three equations, together with \(\widetilde C_k\succ0\), reproduce the three latent moment recursions after setting \(b=m_1-Wm_0\); the remaining stability, logarithm, Stein, and supplied-noise equality restrictions are what distinguish the present structural subclass.

By \(\text{lem:aalto-model-scope}\), the published full-state moment recursion is
\[
 P_{j+1}=AP_jA^{\mathsf T}+Q,
 \qquad
 \bar m_{j+1}=A\bar m_j+b.
\]
Setting \(A=W\), \(P_j=\widetilde C_j\), \(Q=0\), and \(b=(I_p-W)x^\star\) maps every present noise-subtracted equation into that recursion.  The implication is only equation-level: the cited model directly observes its state, whereas the present model observes a Gaussian state plus a supplied common measurement error.  It is also a strict inclusion.  For example, the cited recursion allows the invertible choice \(A=2I_p\), \(Q=0\), and any positive-definite initial covariance; this choice is excluded here because its spectral radius is \(2\), while every \(W=\exp(hG)\) allowed by \(\text{ass:hurwitz}\) has spectral radius below one.  The cited scope further allows known nonzero \(Q\), imposes no principal-logarithm embedding condition, and, by \(\text{lem:aalto-model-scope}\), does not require Gaussian laws once finite second moments exist.  This proves the asserted strict subclass relation and shows why no equality with the published observation-model class follows.

Finally, we derive the finite sign map rather than importing it.  Define the positive-definite whitened matrices and the whitened transition
\[
 R_0=\widetilde C_0^{-1/2}\widetilde C_1\widetilde C_0^{-1/2},
 \qquad
 R_1=\widetilde C_1^{-1/2}\widetilde C_2\widetilde C_1^{-1/2},
 \qquad
 U=\widetilde C_1^{-1/2}W\widetilde C_0^{1/2}.
\]
The first covariance equation is equivalent to \(UU^{\mathsf T}=I_p\), and substitution into the second gives
\[
 UR_0U^{\mathsf T}=R_1.
\]
Suppose that \(R_0\) and \(R_1\) have the same simple ordered spectrum.  Choose orthogonal eigendecompositions
\[
 R_0=P\Lambda P^{\mathsf T},
 \qquad
 R_1=Q_1\Lambda Q_1^{\mathsf T},
 \qquad
 \Lambda=\operatorname{diag}(\lambda_1,\ldots,\lambda_p),
\]
where \(\lambda_i\ne\lambda_j\) for \(i\ne j\).  Then \(D=Q_1^{\mathsf T}UP\) is orthogonal and satisfies \(D\Lambda D^{\mathsf T}=\Lambda\), equivalently \(D\Lambda=\Lambda D\).  Entrywise,
\[
 (\lambda_j-\lambda_i)D_{ij}=0.
\]
Simplicity therefore makes \(D\) diagonal, and orthogonality gives \(D_{ii}\in\{-1,1\}\).  Hence every covariance-compatible transition belongs to the explicit list
\[
 W_D=\widetilde C_1^{1/2}Q_1DP^{\mathsf T}\widetilde C_0^{-1/2},
 \qquad
 D=\operatorname{diag}(s_1,\ldots,s_p),
 \qquad
 s_i\in\{-1,1\}.
\]
There are at most \(2^p\) such candidates.  Put
\[
 a=P^{\mathsf T}\widetilde C_0^{-1/2}d_0,
 \qquad
 c=Q_1^{\mathsf T}\widetilde C_1^{-1/2}d_1.
\]
The mean equation \(Wd_0=d_1\) is equivalent to \(Da=c\).  If every \(a_i\ne0\), each sign is forced by the computable ratio
\[
 s_i=\frac{c_i}{a_i},
\]
and compatibility requires every such ratio to belong to \(\{-1,1\}\).  For four-snapshot compatibility one additionally tests
\[
 W_D\widetilde C_2W_D^{\mathsf T}=\widetilde C_3,
 \qquad
 W_Dd_1=d_2.
\]
Testing the survivors also against \(\rho(W_D)<1\), the negative-real spectral cut, the Stein equation, and \(E_{W_D}=E\) gives exactly the members of \(\mathcal F_{\mathrm{known}}(E)\), because this list exhausts the first two covariance congruences and every remaining test is one of the defining restrictions.  This completes both the exact known-noise reduction within the present model and the finite sign classification, without asserting equality with the broader published class.
